\clearpage\fancyhead[L]{\small COMET.}\hypertarget{native-342-lower}{}\label{native-342-lower}
\begin{center}\textbf{COMET.}\par\vspace{4pt}\hrule\vspace{4pt}\fontsize{10}{12}\selectfont Göttingische gelehrte Anzeigen. Stück 151. Seite 1497\,..\,1499. 1811 Sept. 21.\par\vspace{4pt}\hrule\vspace{4pt}\end{center}

\par\indent Der grosse diesjährige Comet, dessen Wiedererscheinen wir im 130. Stück unserer Gel. Anz. [S. 327 d. B.] im Voraus angekündigt hatten, wurde hier von unsern Astronomen zum ersten Male den 22. August tief in der Abenddämmerung gesehen, aber eigentliche Beobachtungen konnten erst später angestellt werden, da die Aussicht von der Sternwarte auf der Nordseite durch die Gebäude der Stadt zu sehr beschränkt ist. Es wurden Distanzen des Cometen von Fixsternen gemessen, welche Methode bei diesem hellen Cometen sehr wohl anwendbar ist, und besonders wenn die Beobachtungen vom Stativ gemacht werden, viele Genauigkeit gibt. Bis jetzt (10. September) sind indess von den auf diese Weise angestellten Beobachtungen nur erst folgende des Hrn. Prof. \textsc{Gauss} vollständig reducirt:

\Needspace{9\baselineskip}
\begin{center}\fontsize{10}{12}\selectfont\setlength{\tabcolsep}{4pt}\renewcommand{\arraystretch}{1.15}\begin{tabular}{ll|c|c|c}
\multicolumn{2}{c|}{1811} & Mittlere Zeit & ger. Aufsteigung & Nordl. Abweich. \\\hline
Sept. & 4 & $8^{\mathrm u}\,28^{\mathrm m}\,47^{\mathrm s}$ & $158^\circ\;25'\;24''$ & $39^\circ\;18'\;2''$ \\
 & 6 & $8\;\;48\;\;38$ & $160\;\;\;23\;\;\;16$ & $40\;\;\;14\;\;\;16$ \\
 & 7 & $8\;\;57\;\;\;6$ & $161\;\;\;26\;\;\;10$ & $40\;\;\;41\;\;\;54$ \\
\end{tabular}
\end{center}

\par\indent Hr. Prof. \textsc{Gauss} hat nach diesen Beobachtungen seine vorläufigen, am angef. O. mitgetheilten, parabolischen Elemente verbessert und folgende Resultate gefunden:

\Needspace{13\baselineskip}
\begin{center}\fontsize{10}{12}\selectfont\setlength{\tabcolsep}{4pt}\renewcommand{\arraystretch}{1.15}\begin{tabular}{l}
Durchgang durch die Sonnennähe 1811 Sept. 12.\\
\quad $5^{\mathrm u}\,21^{\mathrm m}\,15^{\mathrm s}$ m. Z. in Göttingen\\
\end{tabular}\end{center}
\begin{center}\fontsize{9}{11}\selectfont\setlength{\tabcolsep}{2pt}\renewcommand{\arraystretch}{1.15}\begin{tabular}{llr}
Länge der Sonnennähe & .\;.\;.\;.\;.\;.\;.\;.\;.\;.\;.\;.\;.\;.\;.\;. & $75^\circ\,17'\,34''$\\
Kleinster Abstand von der Sonne & .\;.\;.\;.\;.\;.\;.\;.\;.\;.\;.\;. & $1.04006$\\
Länge des aufsteigenden Knoten & .\;.\;.\;.\;.\;.\;.\;.\;.\;.\;.\;. & $140^\circ\,24'\,13''$\\
Neigung der Bahn & .\;.\;.\;.\;.\;.\;.\;.\;.\;.\;.\;.\;.\;.\;.\;.\;. & $73\;7\;16$\\
\end{tabular}\end{center}
\noindent Bewegung \emph{rückläufig}.

\par\indent Von einer Ellipticität der Bahn ist bisher noch keine Spur zu erkennen.

\par\indent Um den Lauf des Cometen während der nächsten Monate desto besser übersehen zu können, hat Hr. Prof. \textsc{Gauss} nach obigen verbesserten Elementen eine kleine Ephemeride berechnet, woraus erhellt, dass der Comet den 15. October der Erde am nächsten kömmt, aber dann noch fast um den vierten Theil weiter von ihr absteht, als die Sonne. Seine grösste Lichtstärke erreicht er, falls er überhaupt nur durch\SourceLeaf{343} entlehntes Licht sichtbar ist, am 7. October, und er ist dann etwa fünfzehn Mal so hell, als er am 11. April war. Den 3. October kommt der Comet dem Stern $\eta$ im grossen Bär, und den 3. December dem Stern $\alpha$ im Adler bis auf weniger als Einen Grad nahe.

\Needspace{15\baselineskip}\noindent Lauf des Cometen während der letzten Monate dieses Jahres:

\begin{center}\fontsize{10}{12}\selectfont\setlength{\tabcolsep}{3pt}\renewcommand{\arraystretch}{1.15}\begin{tabular}{ll|c|c|c|c}
\multicolumn{2}{c|}{1811} & \shortstack{Gerade\\Aufsteigung} & \shortstack{Nordliche\\Abweichung} & \shortstack{Abstand\\v. d. Erde} & \shortstack{Licht-\\stärke} \\\hline
Sept. & 12\ $5^{\mathrm u}$ & $166^\circ\;54'$ & $42^\circ\;58'$ & 1.614 & 0.355 \\
 & 22\ $6$ & $182\;\;\;2$ & $47\;\;\;17$ & 1.443 & 0.432 \\
Octob. & 2\ $7$ & $203\;\;\;24$ & $49\;\;\;34$ & 1.306 & 0.490 \\
 & 12\ $8$ & $228\;\;\;45$ & $47\;\;\;16$ & 1.233 & 0.492 \\
\end{tabular}
\end{center}

\Needspace{15\baselineskip}\noindent Lauf des Cometen während der letzten Monate dieses Jahres:

\begin{center}\fontsize{10}{12}\selectfont\setlength{\tabcolsep}{3pt}\renewcommand{\arraystretch}{1.15}\begin{tabular}{ll|c|c|c|c}
\multicolumn{2}{c|}{1811} & \shortstack{Gerade\\Aufsteigung} & \shortstack{Nördliche\\Abweichung} & \shortstack{Abstand\\v. d. Erde} & \shortstack{Licht-\\stärke} \\\hline
Octob. & 22\ $0^{\mathrm u}$ & $251^\circ\;23'$ & $39^\circ\;56'$ & 1.248 & 0.421 \\
Nov. & 1\ $10$ & $267\;\;\;58$ & $30\;\;\;32$ & 1.353 & 0.310 \\
 & 11\ $11$ & $279\;\;\;38$ & $22\;\;\;18$ & 1.531 & 0.208 \\
 & 21\ $12$ & $288\;\;\;13$ & $14\;\;\;55$ & 1.744 & 0.138 \\
Dec. & 1\ $13$ & $294\;\;\;54$ & $9\;\;\;42$ & 1.984 & 0.093 \\
 & 11\ $14$ & $300\;\;\;24$ & $5\;\;\;49$ & 2.232 & 0.064 \\
 & 21\ $15$ & $303\;\;\;22$ & $3\;\;\;0$ & 2.426 & 0.048 \\
\end{tabular}
\end{center}

{\rightskip=0pt plus 1em\par\indent Sonderbar ist die Form des \emph{Schweifes}, der sich um den Cometen herumbiegt, ohne mit ihm selbst zusammen zu hängen, und in zwei gegen einander geneigte Aeste ausläuft. Von einem eigentlichen \emph{Kerne} ist übrigens in dem Cometen gar keine Spur zu sehen.\par}
